\section{Proof of Concept}
\label{sec:proof-of-concept}

\subsection{Scope and Experimental Setup}
\label{subsec:poc-setup}

The service-level intent-assurance loop of Section~4 is instantiated in an
Open5GS/UERANSIM testbed with three simultaneously active service classes. The
scenario resembles a small operational network in which not all users have the
same service contract. One UE represents a premium, latency-sensitive service;
a second UE represents a protected control service that should not be harmed by
remediation; and a third UE represents best-effort background traffic. The
premium service is the target of the intent. Its SLA requires high bidirectional
throughput, sub-$10$~ms latency, low jitter, and negligible packet loss. The
control service is monitored as a guardrail, while the background service is
allowed to receive weaker treatment.

Each UE establishes a real 5G data session through the emulated RAN and the
Open5GS core. Application traffic is generated against a dedicated
\texttt{iperf3} endpoint for each service class, so the measurements reflect
the path through the UE tunnel, gNB, AMF/SMF session control, and UPF
forwarding policy. The telemetry available to the agent combines service-level
measurements and network state: rolling throughput, latency, jitter, and loss
windows; UE registration and tunnel status; network-function logs; subscriber
records; and relevant Open5GS configuration. The agent can also query these
signals repeatedly after taking an action, allowing the loop to distinguish
between a repair that merely changes configuration and a repair that actually
restores service.

The agent is instructed to diagnose before acting. Its control surface mirrors
the operations that would be available to a slice operator: inspect subscriber
profiles, inspect AMF/SMF and UPF configuration, measure the affected UE,
compare the premium service against the control service, update the UPF QoS
profile, repair subscriber and slice mappings, restart a network function when
required, and trigger UE re-registration when a subscriber change must take
effect. The important point is that the agent is not given the injected fault.
It receives the live intent, observes symptoms, selects diagnostic checks, and
then chooses the smallest corrective action that is consistent with the
evidence.

Service degradation is introduced through three representative fault classes.
The first demotes the premium service at the UPF, causing premium traffic to be
treated like background traffic. The second injects packet loss into the
premium class while leaving the nominal slice assignment unchanged. The third
alters subscriber state so that the premium UE is effectively associated with
the wrong service class. These faults exercise different parts of the
assurance loop. UPF faults are visible as degraded user-plane measurements;
subscriber faults require correlating service symptoms with control-plane and
subscriber state. Compliance is computed over the final stabilization window
from the formal intent objectives, with separate score components for premium
throughput, latency, jitter, packet loss, and protected-control impact.

\subsection{Results}
\label{subsec:poc-results}

\begin{figure}[t]
    \centering
    \includegraphics[width=\linewidth]{logs/analysis/figures/comparison.pdf}
    \caption{Per-metric SLA compliance over five repetitions for each fault
    condition and model ($N=40$ trials). Each stacked segment contributes one
    sixth of the score; whiskers show the min--max range across trials.}
    \label{fig:poc-comparison}
\end{figure}

\begin{figure}[t]
    \centering
    \includegraphics[width=\linewidth]{logs/analysis/figures/narrative.pdf}
    \caption{Temporal recovery trace for a demoted premium slice.
    The premium latency initially violates the $10$~ms intent, the agent
    inspects and measures the service, applies the UPF QoS repair, and the
    final recovery window remains compliant.}
    \label{fig:poc-recovery-trace}
\end{figure}

The aggregate compliance results in Fig.~\ref{fig:poc-comparison} indicate
that SLA restoration is reliable when the impairment is expressed in the UPF
service-class configuration. \texttt{gpt-5.4} attains full compliance across
all five baseline, demoted, and lossy repetitions, with mean scores of $1.00$
in each case. The result is significant because the demoted and lossy cases do
not simply ask the agent to recognize an alarm; they require it to connect an
SLA symptom to the data-plane mechanism that can explain it. In successful UPF
repairs, the agent first verifies that the premium service is degraded, checks
that the reference service remains stable, inspects or infers the active QoS
treatment, and then restores the premium profile. The closed loop is completed
only after the post-action measurements satisfy the original intent.

The only residual failure for \texttt{gpt-5.4} occurs under subscriber
misrouting, where one trial ends without restored connectivity and the mean
score decreases to $0.80$. This fault is harder because the visible symptom can
look like a service-quality problem while the cause is stored in subscriber
state and may require UE re-registration to become effective. The agent must
therefore reason across layers: a wrong subscriber association can produce
missing or degraded connectivity even when the UPF policy itself is healthy.

\texttt{gpt-5.4-mini} maintains full compliance in the healthy baseline and
subscriber-misrouting runs, but exhibits lower robustness under impaired
user-plane conditions. Its mean score is $0.80$ for service-class demotion and
$0.77$ for lossy premium traffic. The failed score components are concentrated
in premium latency, jitter, and loss, indicating incomplete or delayed
restoration of the target traffic class rather than a broad control-plane
failure. The comparison therefore separates two dimensions of capability:
whether the agent can eventually identify the correct subsystem, and whether it
can do so consistently within the measurement and stabilization window.

The temporal trace in Fig.~\ref{fig:poc-recovery-trace} illustrates the same
behavior for a single demoted-slice trial. Before remediation, premium latency
is close to $35$~ms and exceeds the $10$~ms SLA reference, while the control UE
remains near its expected low-latency range. This contrast gives the agent a
useful diagnostic cue: the problem is specific to the premium treatment rather
than a general collapse of the testbed. The corrective action restores the UPF
QoS profile associated with the premium class. After the update, premium
latency falls to approximately $2$--$5$~ms and remains below the SLA threshold
throughout the recovery window. The absence of a network-function restart
indicates that the violation can be corrected at the user-plane policy layer
without disrupting the protected control slice. In operational terms, the
agent performs the desired assurance behavior: it observes an SLA violation,
localizes the likely cause, applies a bounded repair, and verifies that the
declared service intent is again satisfied.
